Estimated time of arrival (ETA) predicts trip duration under evolving traffic and underpins navigation, dispatch, and logistics; small errors degrade routing, coordination, and user trust.

Machine learning approaches to ETA require datasets that align road topology, vehicle motion, and labels in a form models can learn from.
We present a toolchain that produces \emph{rich, feature-heavy dynamic graph} data suitable for models that use explicit graph structure, for example architectures with separate treatment of static network elements and time-varying traffic.

The system supports two complementary paths.
First, \emph{simulation-based} datasets: the user supplies a SUMO network and simulation configuration; the tool materializes spatio-temporal graph snapshots with detailed per-node and per-edge features from microsimulation.
Second, \emph{trajectory-based} datasets: real vehicle trajectories are imported and aligned to the same graph representation via map matching and step conversion on a user-selected region.
Porto taxi trajectories illustrate real-world conversion in our bundle; other CSVs and regions are supported.

To demonstrate the value of this representation for learning and comparison, we also built a \emph{web-based testing environment} that couples the simulator to interactive trips, model inference, and logged prediction error, so ETA models can be evaluated under controlled, repeatable conditions.

\paragraph{Keywords.}
\begin{sloppypar}
Estimated time of arrival; traffic simulation; SUMO; dynamic graphs; dataset generation; graph neural networks; web-based evaluation.
\end{sloppypar}

\paragraph{Report scope.}
This document presents the toolchain, evaluation methodology, and results in the depth expected for an advanced M.Sc. project report, including user-interface description, implementation context, and appendices. The narrative is self-contained and can be read without reference to any separate publication.
